\chapter[Training Diagnostics]{Training Diagnostics: What Learning Looks Like at
\texorpdfstring{$R^2 = 0.01$}{R2 = 0.01}}
\label{ch:diagnostics}

Most practitioners' instincts about training curves were calibrated on
problems where models reach 90\% accuracy. Those instincts misfire badly at
$R^2 = 0.01$, in both directions: healthy runs look broken, and broken runs
look healthy. This chapter recalibrates the eyes, then gives the debugging
ladder we used when things actually were broken.

\section{Recalibrating: numbers that look wrong but are right}

\begin{itemize}
\item \textbf{The loss barely moves --- forever.} Training on
  $1 - \Rw$, a perfect model converges to $\approx 0.986$; a useless one
  sits at $1.000$. The \emph{entire dynamic range of the problem} lives
  between 1.00 and 0.98. Plot the third and fourth decimals or plot
  nothing. (Our logs print losses to five places for exactly this
  reason.)
\item \textbf{Epoch 1 is worse than doing nothing.} A freshly initialized
  network outputs nonzero noise, which scores \emph{negative} against the
  zero-origin metric (train $R^2 \approx -0.02$ is typical). It should
  cross zero within an epoch or two and creep upward. Panic only if it
  doesn't cross.
\item \textbf{Validation $R^2$ of $+0.008$ is a good model.} Internalize
  the ladder of Chapter~\ref{ch:metric}. And its converse ---
\item \textbf{validation $R^2$ of $+0.05$ is an emergency}, not a
  breakthrough: at 3--4$\times$ the known ceiling, the probability that
  you out-modeled the world's best team by that margin is far below the
  probability of a centered rolling window
  (Chapter~\ref{ch:leakage}). Our standing rule: celebrate \emph{after}
  the leak hunt fails.
\item \textbf{Train--validation gaps are structural.} Train $R^2$ of
  0.02--0.03 against validation 0.008 is the equilibrium of a correctly
  regularized model at this SNR --- memorization pressure is constant and
  nonzero. The gap is monitored for \emph{growth}, not for existence.
\end{itemize}

\section{The curves worth logging}

Beyond the loss, five series earn their disk space in every run:

\begin{enumerate}
\item \textbf{Validation weighted-$R^2$ per epoch} --- the decision
  variable for early stopping (patience $\sim$10: the curve is noisy
  enough that patience 2--3 stops on fluctuations; we measured
  best-epoch spreads of 5+ epochs across seeds of identical configs).
\item \textbf{Per-day validation score} at the end --- a distribution, not
  a number. A model that wins by $+0.02$ on three days and loses
  everywhere else is fragile in a way the mean hides; a model whose
  worst days cluster at regime events tells you what feature is missing.
\item \textbf{Prediction statistics}: mean (bias burns score
  quadratically --- Chapter~\ref{ch:metric}), standard deviation
  (collapse toward zero-variance predictions is the classic
  under-trained/over-regularized signature; explosion signals a broken
  scaler), and max $|\hat y|$ (one 50-sigma prediction on a high-weight
  row can erase a day).
\item \textbf{Gradient norms} pre-clipping --- heavy-tailed days announce
  themselves here first; a rising trend mid-training precedes most
  divergences we saw.
\item \textbf{The static/online pair} for any adaptive model
  (Chapter~\ref{ch:online}) --- the gap between them is a diagnostic
  in its own right, and its collapse is how you notice a refit-lr
  regression before it costs a week.
\end{enumerate}

\section{The debugging ladder}

When a model underperforms expectations, resist touching hyperparameters
--- at this SNR you cannot distinguish a bad lr from a broken pipeline by
staring at scores. Climb in order; each rung is cheap and conclusive:

\subsection{Rung 1: anchor the floor}

Score the trivial ensemble: all-zeros (must be exactly $0.000$ --- if not,
your metric implementation is wrong, and yes, we have seen that), plus a
10-feature ridge probe (should land $+0.002$ to $+0.004$ here). Your model
must live above both. A deep model \emph{below the linear probe} has a
wiring problem, not a capacity problem.

\subsection{Rung 2: overfit on purpose}

Train on 3 days, validate on the same 3 days. Any healthy architecture
must push train $R^2$ far above the signal level ($0.2+$) by memorizing.
Failure localizes the bug precisely: data loader shapes, mask direction,
loss reduction, weight plumbing --- something mechanical. This test found
a transposed (symbol, time) batch for us in minutes; no amount of
production-scale training would have.

\subsection{Rung 3: causality and structure probes}

The shift and shuffle tests of Chapter~\ref{ch:leakage}, run as
\emph{model} diagnostics: score collapse under a one-step feature shift
$\Rightarrow$ leakage; score \emph{survival} under within-day shuffling
$\Rightarrow$ your sequence model is not using sequence (which, for an
RNN scoring near the MLP baseline, is confirmation of rung-2-class bugs
or of features that already contain all the time structure).

\subsection{Rung 4: seed replication}

Three seeds minimum before comparing anything to anything
(Chapter~\ref{ch:problem}). Our measured seed spread ($\pm 0.0003$ to
$0.0008$) is the unit of measurement; an effect inside one unit does not
exist yet.

\subsection{Rung 5: now, and only now, hyperparameters}

And even then, sweep the parameters with cliffs first (refit lr ---
Chapter~\ref{ch:online}; training lr; clipping threshold) before the
parameters with plateaus (width, depth, dropout) --- our capacity sweeps
were flat, our procedure sweeps were not.

\section{Comparing runs honestly}

The bookkeeping that turns runs into knowledge
(Chapter~\ref{ch:validation} covers the validation side):

\begin{itemize}
\item \textbf{One variable per comparison.} The temptation to bundle
  (``new features \emph{and} more width, it won!'') destroys attribution
  exactly where effects are small enough that attribution is the whole
  point.
\item \textbf{Fixed harness, forever.} Same dates, same embargo, same
  metric code, same seeds across every experiment in a campaign. All of
  this book's tables are mutually comparable because the harness never
  moved.
\item \textbf{Append-only results log} with config hashes. Selection bias
  feeds on forgotten losers (Chapter~\ref{ch:leakage}); the log is the
  immune system. Ours is a directory of JSON artifacts plus two markdown
  ledgers --- nothing fancier than that is needed.
\end{itemize}

\begin{jsbox}[: a real debugging transcript, compressed]
Symptom: sig-transformer online score $0.0044$, well below its static
$0.0059$ --- ``online refit breaks signatures.''
Rung 1--2: pass (wiring fine). Rung 3: pass (no leak; sequence used).
Rung 4: reproduces across seeds --- so it is real, but real
\emph{what}? Rung 5, cliff parameters first: refit-lr sweep $\to$ at
$3\times10^{-5}$ the regression becomes a gain ($+0.0062$).
Conclusion rewritten: not an architecture defect --- a borrowed
hyperparameter (Chapter~\ref{ch:online}). Total ladder time: one
afternoon. Time spent on the wrong conclusion before running the ladder:
two weeks of intermittent confusion.
\end{jsbox}

\begin{drillbox}
Instrument your current training script with the five curves of this
chapter, then deliberately break it four ways --- transpose two batch
dimensions, unweight the loss, center a rolling window, set refit lr
10$\times$ high --- and re-run. Learn each failure's curve signature
while the stakes are zero; at $R^2=0.01$, recognizing them from curves is
a survival skill, because the \emph{scores} of broken and healthy runs
overlap.
\end{drillbox}
